\chapter{Testing and Results}
\label{ch:testing}

\section{Testing Strategy}
Testing is a critical phase of the practical training that ensures our system is reliable, secure, and performs correctly. During my training, I used a multi-level testing strategy:
\begin{itemize}
  \item \textbf{Manual Functional Testing:} I manually ran the Flutter mobile app and typed in various scam and safe messages to check if the on-screen alerts, local rule engine, and backend machine learning classifier worked as expected.
  \item \textbf{Backend Verification:} I executed the backend JUnit tests to verify database connections, authentication rules, and user-role permissions.
  \item \textbf{Web Dashboard Build Verification:} I ran Vite type-checking and build commands on the React client to make sure the interface has no TypeScript compilation errors.
  \item \textbf{Security Scanning:} I analyzed the repository's Continuous Integration (CI) configuration, which includes tools like Gitleaks (to scan for exposed passwords/keys) and OWASP ZAP (to run automated security tests on the Spring Boot backend).
\end{itemize}

\section{Test Environment}
I conducted my testing using a local development setup. The backend, PostgreSQL database, and machine learning service were run inside Docker containers using Docker Compose. The Flutter application was compiled and tested using an Android emulator (targeting Android API level 33). The React frontend was tested locally using a standard web browser on port 5173.

\section{Functional Test Cases}
To verify that the system satisfies all system requirements, I designed and executed several functional test cases. The inputs, expected behaviors, and actual observed results of these tests are documented in Table~\ref{tab:functional-tests}.

\begin{table}[htbp]
  \centering
  \caption{Functional testing results of the Argus platform}
  \label{tab:functional-tests}
  \begin{tabularx}{\textwidth}{l p{0.25\textwidth} p{0.30\textwidth} X}
    \toprule
    \textbf{ID} & \textbf{Scenario / Input} & \textbf{Expected Result} & \textbf{Observed Status}\\
    \midrule
    FT-01 & Manually enter a safe text message in Flutter app. & The app displays a green "Safe" badge and threat score of ~0.02. & Passed. Local rules processed input instantly.\\
    FT-02 & Manually enter an M-Pesa scam message in Flutter. & The app displays a red "Fraud" alert and warns the user. & Passed. Local rules triggered and flagged the keywords.\\
    FT-03 & Send message to server with active internet connection. & Spring Boot forwards message to FastAPI BERT; returns AI classification. & Passed. FastAPI predicted label with high confidence.\\
    FT-04 & Log in to administrator dashboard with admin credentials. & React displays dashboard structure and mock statistics panels. & Passed. React rendered the UI without compilation errors.\\
    \bottomrule
  \end{tabularx}
\end{table}

\section{Manual Prediction Examples}
To test the accuracy of our dual-detection approach, I created a set of test SMS messages. These represent real-world scam and safe formats commonly found in Tanzania. I manually submitted these messages to our scanner and logged the system outputs in Table~\ref{tab:example-predictions}.

\begin{table}[htbp]
  \centering
  \caption{Manual scan verification with sample SMS messages}
  \label{tab:example-predictions}
  \begin{tabularx}{\textwidth}{X p{0.15\textwidth} p{0.15\textwidth} p{0.25\textwidth}}
    \toprule
    \textbf{Test Message Body} & \textbf{Source Checked} & \textbf{System Verdict} & \textbf{My Interpretation}\\
    \midrule
    "Congratulations! You have won 1,000,000 TZS in our Vodacom lottery. To claim your prize, please call 0712XXXXXX immediately." & Heuristic Rules \& AI Model & \textbf{Fraud} (Threat: 99\%) & Correctly detected based on words "won", "lottery", and "prize".\\
    "M-Pesa: Umepokea TZS 50,000 kutoka kwa Zephaniah. Salio lako jipya ni TZS 120,000." & Heuristic Rules \& AI Model & \textbf{Safe} (Threat: 1\%) & Correctly classified as a standard transaction confirmation.\\
    "Ndugu mteja, akaunti yako ya benki imefungwa. Tafadhali bofya kiungo hiki ili kuithibitisha: http://fake-bank-login.com" & Heuristic Rules \& AI Model & \textbf{Fraud} (Threat: 95\%) & Correctly flagged due to Swahili bank scam patterns and external link.\\
    "Hi Zephaniah, are we still meeting for lunch at 1 PM today?" & AI Model Only & \textbf{Safe} (Threat: 2\%) & Correctly identified as normal personal communication.\\
    \bottomrule
  \end{tabularx}
\end{table}

\section{Analysis of Heuristics vs. AI Machine Learning}
During my manual testing, I observed several interesting behaviors regarding how the two detection engines work together:
\begin{itemize}
  \item \textbf{Speed:} The local rule-based engine is incredibly fast, returning a threat score in under 5 milliseconds. This is because it runs completely on the device's CPU and does not require any network calls.
  \item \textbf{Context Understanding:} While the rule engine is excellent for exact keywords, it can be fooled by spelling mistakes (e.g., "pr1ze" instead of "prize"). The server-side BERT model, however, was able to correctly flag these messages because it reads the overall context of the sentence instead of looking for exact word matches.
  \item \textbf{Network Fallback:} When I simulated an offline environment by disabling internet access on the Android emulator, the app handled it gracefully. It immediately skipped the backend call, used the local rules, and warned me that the scan was completed offline.
\end{itemize}

\section{Discussion of Testing Achievements}
Through these testing activities, I verified that:
\begin{enumerate}
  \item The background Android SMS listener successfully intercepts incoming messages and fires local alerts.
  \item The Spring Boot API correctly acts as an intermediary, forwarding messages to FastAPI and saving only malicious scam scans in PostgreSQL.
  \item The React client successfully loads its dashboard and handles in-memory changes, although it still needs to be fully integrated with live backend REST services to display dynamic administrative data in a production environment.
\end{enumerate}

\section{Chapter Summary}
In this chapter, I documented my testing strategy, functional test cases, and manual scanning verification with sample messages. I analyzed the speed and intelligence trade-offs between local rules and server-side BERT models, confirming that the fallback mechanisms worked flawlessly. In the next chapter, I will discuss the wider implications, security concerns, and lessons learned during my training.
